% m2_model.tex -- Main text, Section 2: Model, scenes and parameter regime.
% Paths in "% src:" comments are relative to the repository root.
% Numbers of the two tables and of Section 2.6: code/s2_model_scene_table.py ->
% data/s2_model_scenes.json (cross-checked there against data/plan_theory_checks.json
% to 2e-13) and data/s2_model_table_rows.tex (table rows as generated).
\section{Model, scenes and parameter regime}\label{s2_model:sec}

This section fixes the model: one movement rule, a mean-field model with a contact kernel, an individual-based
model of which it is the mean field, and three reproduction numbers that must be kept apart. It then describes
the four scenes and states the physical regime of the parameters. In Sections~\ref{s2_model:sec}--\ref{s6_scenes:sec}
lengths are measured in lattice cells and time in days.

\subsection{Movement on the lattice}\label{s2_model:sec:movement}

The floor of a room is covered by a square lattice of spacing $a$. Cells occupied by walls, furniture or other
barriers are removed; the remaining $\ncell$ walkable cells form the set $\cells$, and two walkable cells are
neighbours if they share a side. Cell $x$ has a mobility (diffusion coefficient) $D_x=\Dz d_x>0$, where $\Dz$
is the mobility scale and $d_x$ a dimensionless zone multiplier. A person in cell $x$ hops to a neighbouring
cell $y$ at rate
\begin{equation}\label{s2_model:eq:hop}
  \hop{x}{y}=\hop{y}{x}=\frac{2D_xD_y}{D_x+D_y}\qquad\text{for neighbouring cells } x,y\in\cells ,
\end{equation}
the harmonic mean of the two mobilities. This is the interface coefficient of finite-volume discretisations
of diffusion in composite media \citep{patankar1980numerical}: the net flow between two cells is
$\hop{x}{y}$ times the difference of their occupation numbers, so that \eqref{s2_model:eq:hop} is a
discretisation of Fickian diffusion, $\partial_t\rho=\nabla\!\cdot(D\nabla\rho)$, with reflecting walls. The
movement generator is the $\ncell\times\ncell$ matrix
\begin{equation}\label{s2_model:eq:generator}
  M_{yx}=\hop{x}{y}\ (y\neq x),\qquad M_{xx}=-\sum_{y\neq x}\hop{x}{y},\qquad \one^{\T}\gen=0 ,
\end{equation}
with $M_{yx}=0$ if $x\ne y$ are not neighbours. The column index is the source cell: the vector of expected
occupation numbers obeys $\dot N=\gen N$, and $\one^{\T}\gen=0$ says that nobody enters or leaves through the
walls. Because $D_x=\Dz d_x$, we have $\gen=\Dz\genz$, where $\genz$ is the generator at $\Dz=1$. The
propagator $P_t(y\,|\,x)=\bigl(\e^{\gen t}\bigr)_{yx}$ is the probability that a walker that started in cell
$x$ is in cell $y$ at time $t$. Closed forms of the propagator are known for homogeneous rectangular lattices
\citep{giuggioli2020exact}; for the heterogeneous scenes of this article every quantity is computed
numerically from $\gen$.

\begin{assumption}\label{s2_model:ass:A}
  The walkable region is connected: $\gen$ is irreducible.  The population is at movement equilibrium,
  $N_x(t)\equiv N^*_x=\Npop\,\pi_x$, where $\stat>0$ is the stationary distribution of $\gen$
  ($\gen\stat=0$, $\one^{\T}\stat=1$; $\pi_x=1/\ncell$ for the symmetric rule \eqref{s2_model:eq:hop}).
\end{assumption}

Under rule \eqref{s2_model:eq:hop} the matrix $\gen$ is symmetric, so the uniform distribution is the
stationary one and people do not accumulate in slow zones. The assumption is phrased with a general $\stat$
because several results of Section~\ref{s3_r0:sec} hold for any irreducible generator. The harmonic mean is a
modelling choice. A rule in which the rate depends on the departure cell only, $\hop{x}{y}=D_x$, is a different
model: it satisfies detailed balance with $\pi_x\propto1/D_x$, so that people accumulate where they move
slowly, and its continuum limit is $\partial_t\rho=\nabla^2(D\rho)$. With lengths in cells, $D_x$ is in
\cdunit{}; in physical units the rate \eqref{s2_model:eq:hop} is $2D_xD_y/((D_x+D_y)a^2)$.

\subsection{Mean-field model}\label{s2_model:sec:meanfield}

An infective in cell $x$ makes infectious contacts at rate $\beta q_x$, where $\beta$ is the transmission rate
and $q_x\ge0$ the dimensionless infection efficiency of the cell ($\Qm=\diag q$, $q\not\equiv0$). A contact
lands in cell $y$ with probability $P_{xy}=P(y\,|\,x)$, where the contact kernel $\ker$ is row-stochastic, and
reaches a person chosen uniformly among those present in $y$ (frequency-dependent incidence). Infectives
recover at rate $\gamma$. With $S_y,I_y,R_y$ the expected numbers of susceptible, infective and recovered
people in cell $y$,
\begin{equation}\label{s2_model:eq:meanfield}
\begin{aligned}
  \dot S_y &= (\gen S)_y-\frac{S_y}{N^*_y}\,\beta\sum_{x}q_xP_{xy}I_x ,\qquad
  \dot I_y = (\gen I)_y+\frac{S_y}{N^*_y}\,\beta\sum_{x}q_xP_{xy}I_x-\gamma I_y ,\\
  \dot R_y &= (\gen R)_y+\gamma I_y .
\end{aligned}
\end{equation}
The \emph{same-cell kernel} $\ker=\Id$ gives the local incidence $\beta q_yS_yI_y/N^*_y$; the
system is then a susceptible--infective--recovered patch model with one patch per cell, and the results
known for such models apply (Section~\ref{s3_r0:sec}). The simulations use the \emph{1\,m kernel}: $P(y\,|\,x)$
is uniform over the walkable cells whose centres lie within 1\,m of the centre of $x$ (line of sight through
barriers is not checked). It coincides with the same-cell kernel when $a>1$\,m; without it the lattice
spacing would silently set the contact range, with cell areas from $0.25\,\mathrm{m}^2$ to $4\,\mathrm{m}^2$
in the scenes below. The range of 1\,m is a modelling choice.
% src: data/s2_model_scenes.json -> scenes.<name>.D0_cell2_per_day_in_m2_per_day (a^2 = 2.25, 4, 1, 0.25 m^2)

Summing the three equations gives $\dot N=\gen N$ for $N=S+I+R$, so that $N_y(t)=N^*_y$ for all $t$ under
Assumption~\ref{s2_model:ass:A}. Writing $S_y/N^*_y=1-(I_y+R_y)/N^*_y$,
\begin{equation}\label{s2_model:eq:linear}
  \dot I=\Jm I+O\bigl(|I|(|I|+|R|)\bigr),\qquad \Jm=\gen+\Fm-\gamma\Id,\qquad \Fm=\beta\,\ker^{\T}\Qm ,
\end{equation}
with $\Fm=\beta\Qm$ for the same-cell kernel. Two features of \eqref{s2_model:eq:linear} are used repeatedly.
First, $\Jm$ contains neither the number of people nor their density: under frequency-dependent incidence,
crowding does not enter the linearised dynamics. Second, the neglected term is
$-\diag\bigl((I+R)/N^*\bigr)\Fm I\le0$, so that every solution of \eqref{s2_model:eq:meanfield} satisfies
$\dot I\le\Jm I$ componentwise for all $t\ge0$.

\subsection{Individual-based model}\label{s2_model:sec:ibm}

The individual-based model has $\Npop$ people who perform independent continuous-time random walks on
$\cells$ with generator $\gen$ (rule \eqref{s2_model:eq:hop}, uniform $\stat$) and who are placed
independently and uniformly at time $0$. Each person is susceptible, infective or recovered. While an
infective is located at $x$ and a susceptible at $y$, the infective infects that susceptible at rate
\begin{equation}\label{s2_model:eq:pairhazard}
  h(x,y)=\frac{\beta\,q_x\,P_{xy}}{\rhobar},\qquad \rhobar=\frac{\Npop-1}{\ncell}
  \quad(\text{frequency-dependent calibration}),
\end{equation}
where $\rhobar$ is the mean number of \emph{other} people per cell. The other $\Npop-1$ walkers are
independent and uniformly distributed, so the expected rate at which an infective at $x$ infects in a fully
susceptible room at movement equilibrium is $\sum_{y}h(x,y)\,\rhobar=\beta q_x$, exactly the rate assumed in
\eqref{s2_model:eq:meanfield}, for every $\Npop$. In the density-dependent (DD) calibration, used only in
Section~\ref{s6_scenes:sec:interventions} and always labelled DD, $\rhobar$ in \eqref{s2_model:eq:pairhazard}
is frozen at its value $\rhobar_{\mathrm{ref}}$ for a reference occupancy; the hazard of a pair is then fixed
and the contact rate becomes $\beta q_x\rhobar/\rhobar_{\mathrm{ref}}$, proportional to occupancy. Infectives
recover at rate $\gamma$: the model is of SIR type with an exponential infectious period of mean $1/\gamma$
and no latent class.

The process is a continuous-time Markov chain with three kinds of event --- a hop, a recovery, and the
infection of a susceptible located at $y$, at rate $\sum h(x_i,y)$ summed over the current positions $x_i$ of
the infectives --- and is simulated exactly by the direct method of
\citet{gillespie1976general,gillespie1977exact}; there is no time step and hence no discretisation error. The
infection tree (infector, generation, time and cell of every infection) is recorded, so that secondary cases
are counted, not fitted. The algorithm and the verification of the simulator are in Supplementary
Section~\ref{S-appA_numerics:sec}.

Model \eqref{s2_model:eq:meanfield} is the mean-field approximation of this process: it replaces the numbers
of susceptibles and infectives in each cell by their expectations and neglects their correlations. The two
models agree on the infection pressure of a case: a case born at $x$ exerts the expected pressure
$\beta\,\Exp_x\!\int_0^\tau q(X_t)\,\dd t$ over its infectious period $\tau$, which is the column sum
$\noff(x)$ of the next-generation matrix of Section~\ref{s3_r0:sec}. What the mean field ignores is that a
discrete person can be infected only once and that an infective keeps meeting the same few neighbours
(Section~\ref{s5_pair:sec}). The approximation becomes accurate when many people share a contact
neighbourhood and mobility is high. In the office scene with 1\pct{} of the people initially infective, the
largest difference between the mean simulated prevalence curve and the solution of
\eqref{s2_model:eq:meanfield}, both as fractions of $\Npop$, is $0.42$ for $\Npop=100$ and $\Dz=1$ (200 runs),
$0.058$ for $\Npop=1000$ and $\Dz=10$ (200 runs), and $0.004$ for $\Npop=5000$ and $\Dz=100$ (40 runs; checked
only by re-running the same code).
% src: data/bsc_sim/01_validation.json -> F_meanfield_limit (max_abs_diff_curve 0.4226, 0.0585, 0.0037; n_rep 200, 200, 40; I0 = 1, 10, 50)
% src: notes/VERIFICATION.md -> Section 4.2, simulations, summary ("Not re-simulated with the second simulator": test F with N = 5000)
The normalisation \eqref{s2_model:eq:pairhazard} by the mean number of \emph{other} people matters. A rule that
divides instead by the instantaneous number of all people in the cell, the infective included, defines a
different process in which an infective alone in its cell transmits to nobody; at the occupancies of the
scenes it is sub-critical (Supplementary Remark~\ref{S-s5_pair:rem:encounter}).

\subsection{Three reproduction numbers}\label{s2_model:sec:numbers}

Three quantities are called reproduction numbers in this article.
\begin{enumerate}
  \item The \emph{basic reproduction number} $\Rz=\srad(\ngm)$ of the mean-field model, the spectral radius of
    the next-generation matrix $\ngm=\Fm(\gamma\Id-\gen)^{-1}$ (Theorem~\ref{s3_r0:thm:ngm}). It is the threshold
    of \eqref{s2_model:eq:linear}: $\Rz>1$ if and only if the growth rate $\lamone=\sabs(\Jm)$ is positive.
  \item The \emph{pair-level} numbers of the individual-based model (Definition~\ref{s5_pair:def:R1}):
    $\Rone(x)$, the expected number of secondary cases of a lone index case born at $x$ when the other
    $\Npop-1$ people are susceptible and only the index transmits; its average $\Ronebar$ over a uniformly
    placed index; and the spectral radius $\srad(\ngmone)$ of a pair-level next-generation matrix. $\Rone(x)$ is
    exact for the first generation. None of the three is a threshold quantity
    (Remark~\ref{s5_pair:rem:threshold}).
  \item The \emph{counted} number $\Rhat$: the mean number of people infected directly by the index case in
    simulation, read from the infection tree. It depends on the law of the index cell (uniform, or weighted
    by the Perron vector of $\ngm$) and on whether later cases also transmit and so compete for susceptibles;
    both are stated wherever $\Rhat$ is quoted.
\end{enumerate}
In the office scene at $\Dz=1$ they are $\Rz=5.107$, $\Ronebar=2.266$ and $\Rhat=2.265\pm0.004$ (standard
error; 400{,}000 runs, uniform index, only the index transmits).
% src: data/s5_pair_large_sample.json -> office|D0=1 (R0_ngm 5.1073, R1bar_exact 2.2661, mean 2.2648, se 0.0039, n_runs 400000)

\subsection{The four scenes}\label{s2_model:sec:scenes}

\begin{table}[tbp]
\centering
\caption{The four scenes: lattice size in cells, spacing $a$, number $\ncell$ of walkable cells, number
$\Npop$ of people, mean number $\rhobar=(\Npop-1)/\ncell$ of other people per cell, means over the walkable
cells of the mobility multiplier $d=D/\Dz$ and of the infection efficiency $q$, the well-mixed value
$\beta\langle q\rangle/\gamma$ and the upper bound $\beta\qmax/\gamma$ of $\Rz$ for $\beta=0.5\perday$ and
$\gamma=0.14\perday$, and the mean number of cells in the 1\,m contact kernel (1.0 means the same-cell
kernel). The zone values of $d$ and $q$ are in Supplementary Table~\ref{S-s2_model:tab:zones}; all of them are
illustrative assumptions, and none is a measurement.}
\label{s2_model:tab:scenes}
\small
% src: data/s2_model_scenes.json -> scenes.<name> (rows as generated in data/s2_model_table_rows.tex);
%      scene definitions code/bsc_sim/scenes/{office,supermarket,classroom,metro}.json
\begin{tabular}{@{}lcrrrcccccc@{}}
\toprule
Scene & lattice & $a$ (m) & $\ncell$ & $\Npop$ & $\rhobar$ & $\langle d\rangle$ & $\langle q\rangle$
  & $\beta\langle q\rangle/\gamma$ & $\beta\qmax/\gamma$ & kernel cells \\
\midrule
Office & $20\times13$ & 1.5 & 234 & 100 & 0.423 & 0.561 & 1.300 & 4.643 & 5.357 & 1.0 \\
Supermarket & $25\times15$ & 2 & 278 & 150 & 0.536 & 1.239 & 0.914 & 3.263 & 7.143 & 1.0 \\
Classroom & $15\times10$ & 1 & 148 & 80 & 0.534 & 0.192 & 1.318 & 4.706 & 7.143 & 4.6 \\
Metro carriage & $45\times6$ & 0.5 & 270 & 310 & 1.144 & 0.028 & 2.536 & 9.056 & 10.000 & 11.1 \\
\bottomrule
\end{tabular}
\end{table}

\begin{figure}[tbp]
\centering
\includegraphics[width=0.84\textwidth]{fig3_scene_maps_en}
\caption{Floor plans of the four scenes. Each small square is one lattice cell of side $a$; colours mark
the zones, whose mobility multiplier $D/\Dz$ and infection efficiency $q$ are given in the legends; dark
cells are barriers and are not part of $\cells$. The panel titles give the lattice size in cells, the
spacing $a$ and the number of people $\Npop$. The arrangement of the zones and the zone values are
illustrative assumptions (see text).}
% src: figures/fig3_scene_maps_{en,zh}.pdf, script code/bsc_sim/experiments/20_figures.py (fig_scenes), scenes code/bsc_sim/scenes/*.json
\label{s2_model:fig:scenes}
\end{figure}

The model is evaluated on four scenes: an open-plan office, a supermarket, a classroom and a metro carriage
(Table~\ref{s2_model:tab:scenes}, Fig.~\ref{s2_model:fig:scenes}). The same scene files are used for the
theory and for the simulations. The lattice sizes and spacings, the zones with their values of $d$ and $q$,
the numbers of people and the arrangement of the zones on the floor are illustrative assumptions; they are not
measurements and are not derived from any standard. In the office the zones cover 60, 15, 10, 5 and 10\,\% of
the 260 cells of the floor (workstations, corridor, meeting room, kitchen and barriers; the 234 walkable cells
exclude the barriers); in the supermarket shelf racks and
counters act as barriers; the classroom desks are split into a window block ($q=1.2$) and an inner block
($q=2.0$); and the carriage has doors, seats and grab rails.
% src: code/bsc_sim/scenes/office.json (zone_counts 156, 39, 26, 13, 26 of 260 cells); code/bsc_sim/scenes/supermarket.json, classroom.json, metro.json
The metro mobilities follow a crowding rule $d=c\,n_{\mathrm{d}}^{-\eta}$ with $\eta=2$ in the standing area
and $\eta=1$ elsewhere, evaluated at the crowd density $n_{\mathrm{d}}=310/67.5\,\mathrm{m}^2=4.59$ people per
$\mathrm{m}^2$; their area
mean is $0.028$.
% src: data/s2_model_scenes.json -> scenes.metro (density_used_for_metro_rule_per_m2 4.5926, d_mean 0.0282, zones); code/bsc_sim/scenes/metro.json (D_coef, D_exp)
The quantity that enters the bounds of Section~\ref{s3_r0:sec} is the mean of $q$ over walkable cells, which
is $1.300$, $0.914$, $1.318$ and $2.536$. The 1\,m kernel is the same-cell kernel in the office and the
supermarket; it contains the cell and up to four neighbours in the classroom ($4.6$ cells on average) and up
to 13 cells in the carriage ($11.1$ on average). Results for the same-cell kernel in these two scenes are
reported alongside as a sensitivity case.
% src: data/s2_model_scenes.json -> scenes.<name>.q_mean; scenes.<name>.kernel1m_cells_mean (1.0, 1.0, 4.62, 11.13), kernel1m_cells_max (1, 1, 5, 13)

\subsection{Parameters and physical regime}\label{s2_model:sec:regime}

The default epidemiological parameters are $\beta=0.5\perday$ and $\gamma=0.14\perday$, a mean infectious
period of $7.1$ days and $\beta/\gamma=3.571$. They are illustrative values and are not calibrated to any
disease or dataset. The same holds for the default mobility scale $\Dz=1\cellday$, and this value matters
for every result below. It is used as a reference low-mobility value at which the spatial effects of the model
are visible; most results are also given at $\Dz=10$ and
$100\cellday$, and $\Rz$ of every scene at walking mobility (Table~\ref{s2_model:tab:regime}).
% src: data/s2_model_scenes.json -> beta, gamma, beta_over_gamma (3.5714), office_regime.infectious_period_days (7.14)

\paragraph{What $\Dz=1$ means.} In the office ($a=1.5$\,m), $\Dz=1\cellday$ is $2.25\,\mathrm{m}^2$ per day.
The root-mean-square displacement $\sqrt{4Dt}$ of a walker with $d=1$ is then $3$\,m in a day and $8$\,m in a
mean infectious period, in a room of $30\,\mathrm{m}\times19.5\,\mathrm{m}$. A walker in the meeting room of
this scene (26 cells) leaves it at the principal rate $\muZ=0.00877\perday$: its residence time $1/\muZ$ is
$114$ days, sixteen infectious periods. In the metro carriage at $\Dz=1$ a passenger makes on average $0.075$
hops of $0.5$\,m per day: the carriage is frozen on the time scale of an infectious period.
% src: data/s2_model_scenes.json -> office_regime (D0_m2_per_day 2.25, rms_m_per_day_D0 3.0, rms_m_per_infectious_period_D0 8.02), scenes.office.floor_m
% src: data/s2_model_scenes.json -> office_regime (meeting_room_mu_Z_per_day 0.008769, meeting_room_1_over_mu_Z_days 114.04)
% src: data/s2_model_scenes.json -> scenes.metro.hop_rate_mean_per_day_D0=1 (0.0749)

\paragraph{Walking mobility.} For a walker whose velocity autocorrelation decays as
$\langle\mathbf{v}(0)\cdot\mathbf{v}(t)\rangle=v^2\e^{-t/\tau}$, the diffusion coefficient in two dimensions
is $D=\tfrac12\int_0^\infty\langle\mathbf{v}(0)\cdot\mathbf{v}(t)\rangle\,\dd t=v^2\tau/2$
\citep{taylor1922diffusion,kubo1966fluctuation}. With an assumed walking speed $v=1.3\,\mathrm{m\,s^{-1}}$, a
typical value for pedestrians \citep{weidmann1993transporttechnik}, and an assumed persistence time
$\tau=1.5$\,s, this gives $D=1.27\,\mathrm{m^2\,s^{-1}}$, that is, $4.9\times10^{4}\cellday$ for $a=1.5$\,m,
and $2.4\times10^{3}\cellday$ if people walk 5\,\% of the time. Over the four scenes these walking values lie
between $1.4\times10^{3}$ and $4.4\times10^{5}\cellday$ (Table~\ref{s2_model:tab:regime}): $\Dz=1$ is three to
almost six orders of magnitude below the mobility of people who walk.
% src: data/s2_model_scenes.json -> walking (D_m2_per_s 1.2675), scenes.office (walking_D0_cell2_per_day 48672, walking_5pct_D0_cell2_per_day 2433.6); same values in data/bsc_theory/worked_examples.json -> E9_physical_D
% src: data/s2_model_scenes.json -> scenes.<name>.walking_5pct_D0_cell2_per_day (smallest: supermarket 1368.9), walking_D0_cell2_per_day (largest: metro 438048); log10: 3.14 to 5.64 (data/misc_checks.json -> mobility_gap)

\begin{table}[tbp]
\centering
\caption{Basic reproduction number $\Rz$ of the mean-field model ($\beta=0.5\perday$,
$\gamma=0.14\perday$) for the mobility scales $\Dz$ (in \cdunit) used in this article and for walking
mobility; in brackets and in the columns headed `excess', the excess of $\Rz$ over the well-mixed value
$\beta\langle q\rangle/\gamma$ in per cent. For the office and the supermarket the 1\,m kernel is the
same-cell kernel; for the classroom and the metro carriage both kernels are given.
The walking values of $\Dz$ are $v^2\tau/2$ with $v=1.3\,\mathrm{m\,s^{-1}}$ and $\tau=1.5$\,s, converted
with the lattice spacing of each scene, for people who walk 5\pct{} of the time and all the time.}
\label{s2_model:tab:regime}
\footnotesize
% src: data/s2_model_scenes.json -> scenes.<name>.R0_1m, scenes.<name>.R0_samecell (rows as generated in
%      data/s2_model_table_rows.tex); D0 = 1, 10, 100 agree with data/plan_theory_checks.json ->
%      scenes.<name>.R0_1m_D0=*, R0_samecell_D0=*
\setlength{\tabcolsep}{3pt}
\begin{tabular}{@{}llcccccccc@{}}
\toprule
 & & & \multicolumn{3}{c}{$\Rz$ (excess, \%) at $\Dz=$}
 & \multicolumn{2}{c}{walking 5\pct} & \multicolumn{2}{c}{walking} \\
\cmidrule(lr){4-6}\cmidrule(lr){7-8}\cmidrule(l){9-10}
Scene & kernel & $\beta\langle q\rangle/\gamma$ & 1 & 10 & 100 & $\Dz$ & excess & $\Dz$ & excess \\
\midrule
Office & same cell & 4.643 & 5.107 (10.0) & 4.676 (0.72) & 4.646 (0.06) & $2.4\times10^{3}$ & 0.0026 & $4.9\times10^{4}$ & 0.00013 \\
Supermarket & same cell & 3.263 & 4.394 (34.6) & 3.391 (3.93) & 3.275 (0.37) & $1.4\times10^{3}$ & 0.0268 & $2.7\times10^{4}$ & 0.00134 \\
Classroom & 1\,m & 4.706 & 5.866 (24.7) & 4.931 (4.79) & 4.731 (0.53) & $5.5\times10^{3}$ & 0.0098 & $1.1\times10^{5}$ & 0.00049 \\
 & same cell &  & 6.182 (31.4) & 4.960 (5.42) & 4.733 (0.57) & $5.5\times10^{3}$ & 0.0105 & $1.1\times10^{5}$ & 0.00053 \\
Metro carriage & 1\,m & 9.056 & 9.241 (2.1) & 9.164 (1.20) & 9.078 (0.25) & $2.2\times10^{4}$ & 0.0013 & $4.4\times10^{5}$ & 0.00006 \\
 & same cell &  & 9.836 (8.6) & 9.308 (2.79) & 9.086 (0.33) & $2.2\times10^{4}$ & 0.0016 & $4.4\times10^{5}$ & 0.00008 \\
\bottomrule
\end{tabular}
\end{table}

\paragraph{Consequence.} At $\Dz=1$ the heterogeneity of $q$ raises $\Rz$ above the well-mixed value
$\beta\langle q\rangle/\gamma$ by $2$ to $35\pct$ (1\,m kernel), depending on the scene. At walking mobility
the excess is below $0.003\pct$ in the office and below $0.03\pct$ in every scene: in the mean-field model a
closed room in which people walk is well mixed, and $\Rz=\beta\langle q\rangle/\gamma$ for all practical
purposes.
% src: data/s2_model_scenes.json -> scenes.<name>.R0_1m.*.excess_over_wellmixed_pct (D0=1: 10.0, 34.6, 24.7, 2.05;
%      walking_5pct: 0.0026, 0.0268, 0.0098, 0.0013; walking: 0.00013, 0.00134, 0.00049, 0.00006);
%      office at D0 = 2400 and 48700 in data/plan_theory_checks.json -> office.mobility_table (0.0027, 0.00013)
All spatial effects in closed rooms reported in Sections~\ref{s3_r0:sec}--\ref{s6_scenes:sec} are large only at
low mobility, and this holds for the pair-level correction of Section~\ref{s5_pair:sec} as well. Leaks through
a localised door are the exception: their effect grows with mobility (Section~\ref{s4_geometry:sec:leaky}).
Results are therefore given for the low default $\Dz=1$ and for $\Dz=10$ and $100$, where the room approaches
the well-mixed limit; the later sections refer back to this paragraph and do not repeat it. A small $\Dz$
could be read as a stand-in for people who stay at a desk or a seat, but a random walk without home positions
is not an adequate model of such people: a slow walker drifts away and does not return, and the stationary
occupancy is uniform over all walkable cells, corridors and aisles included. Section~\ref{s8b_evidence:sec:contacts}
shows the consequence against recorded contacts, at the mobility that matches recorded contact durations: a
free walk calibrated to the amount and duration of contact in six indoor settings gives each person too many
partners in five of the six settings and no group structure in any.

\paragraph{Closed room.} Unless stated otherwise, each scene is a closed room occupied 24 hours a day by the
same $\Npop$ people for the whole epidemic. This is an idealisation for the office and the classroom, and it
is not a physical scenario for the supermarket and the metro carriage, for which the scene files assume visits
of 27 and 20 minutes;
% src: code/bsc_sim/scenes/supermarket.json, metro.json ("dwell")
Section~\ref{s6_scenes:sec:interventions} gives the corresponding numbers per visit and for a room that is
occupied during part of the day.
